% !TeX encoding = UTF-8
\documentclass[variant=guia,cover=institutional,mode=screen]{ua-engineering-v6-article}
% !TeX encoding = UTF-8
\UASetUniversity{Universidad Austral}
\UASetFaculty{Facultad de Ingeniería}
\UASetDepartment{Departamento de Computación}
\UASetCampus{Pilar}
\UASetCareer{Ingeniería Informática}
\UASetCommission{Comisión A}
\UASetProfessor{Equipo docente}
\UASetSemester{Segundo cuatrimestre 2026}
\UASetCourse{Sistemas Distribuidos}
\UASetDate{2026}


\UASetClassNumber{12}
\UASetClassTitle{Chaos Engineering y validación bajo falla}
\UASetDocKind{Guía de ejercicios}

\title{Clase 12 --- Guía de ejercicios}
\author{\UAProfessor}
\date{\UADate}

\begin{document}

\section{Indicaciones generales}
Resuelva cada ejercicio declarando supuestos, modelo de fallas, garantía y trade-off. Cuando corresponda, construya una ejecución verificable y vincule la respuesta con evidencia observable.

\section{Nivel 1: fundamentos y reconocimiento}
\begin{uaexercise}
Explique por qué testing clásico no alcanza para validar propiedades distribuidas bajo falla.
\end{uaexercise}

\begin{uaexercise}
Defina Chaos Engineering y explique por qué no debe confundirse con “romper sistemas al azar”.
\end{uaexercise}

\begin{uaexercise}
Explique qué es un steady state y por qué es indispensable para formular un experimento.
\end{uaexercise}

\begin{uaexercise}
Construya un ejemplo de steady state para un sistema de checkout.
\end{uaexercise}

\begin{uaexercise}
Defina fault injection y explique qué tipos de fallas pueden inyectarse de forma controlada.
\end{uaexercise}


\section{Nivel 2: ejecuciones y fallas}
\begin{uaexercise}
Diseñe un experimento de latencia artificial sobre una dependencia crítica.
\end{uaexercise}

\begin{uaexercise}
Diseñe un experimento de crash de servicio y explique qué propiedad del sistema intentaría validar.
\end{uaexercise}

\begin{uaexercise}
Explique qué diferencias hay entre inyectar una falla de latencia y una de caída total.
\end{uaexercise}

\begin{uaexercise}
Formule una hipótesis experimental correcta para un servicio de pagos.
\end{uaexercise}

\begin{uaexercise}
Explique por qué una hipótesis vaga como “queremos ver si aguanta” no es suficiente.
\end{uaexercise}


\section{Nivel 3: diseño y comparación}
\begin{uaexercise}
Diseñe un experimento que valide si un circuit breaker abre cuando corresponde.
\end{uaexercise}

\begin{uaexercise}
Diseñe un experimento para verificar si retries correctamente acotados no producen tormenta de tráfico.
\end{uaexercise}

\begin{uaexercise}
Defina blast radius y explique por qué debe controlarse explícitamente.
\end{uaexercise}

\begin{uaexercise}
Diseñe una estrategia para limitar el blast radius de un experimento en producción.
\end{uaexercise}

\begin{uaexercise}
Explique qué condiciones mínimas deberían cumplirse antes de correr un experimento de chaos en producción.
\end{uaexercise}


\section{Nivel 4: integración y defensa}
\begin{uaexercise}
Describa un caso donde un experimento sin rollback claro sería irresponsable.
\end{uaexercise}

\begin{uaexercise}
Explique por qué observabilidad y Chaos Engineering están fuertemente conectados.
\end{uaexercise}

\begin{uaexercise}
Diseñe las métricas mínimas necesarias para evaluar un experimento sobre un pipeline de procesamiento de eventos.
\end{uaexercise}

\begin{uaexercise}
Explique qué rol jugarían las trazas en un experimento de degradación por latencia.
\end{uaexercise}

\begin{uaexercise}
Explique por qué un experimento sin instrumentación suficiente produce más ruido que conocimiento.
\end{uaexercise}


\end{document}
